\begin{remark}[Source comparison for the credited re-derivations]\hypertarget{tag-02YS}{}\label{night:P5.7}
\StatementRecord{night:P5.7}{proved}{02YS}{\SeriesRef{night:P5.3}, \SeriesRef{night:P5.4}, \SeriesRef{night:P5.5}, \SeriesRef{night:P5.6}; \cite[Theorems 1.3--1.4, p.~3; Theorem 6.12, pp.~81--82; Proposition 6.16, p.~85; Theorem 7.14, p.~95]{DuEtAl2026v1}; \cite[title, abstract, contents and introduction, pp.~1--3]{DuEtAl2026v2}}
For the fixed standard $(3,4,\infty)$ data, our re-derived tangent
cohomology $(1,1,1,0)$, smooth one-dimensional Kuranishi germ,
automorphism group $\mathbb C^*$, and equivalence
$X_c\simeq X_{c'}\Longleftrightarrow c'-c\in2\mathbb Z$
agree with the corresponding assertions of
\cite{DuEtAl2026v1}. Version~2, entitled
\emph{Oka Geometry of Complex Six-Spheres}, omits those
tangent-cohomology, miniversality, full-automorphism and exact-period
assertions. The tangent and deformation calculations in this appendix
were re-derived after reading version~1. Our working-edition chronology records the public proof of the exact
$2\mathbb Z$ equivalence by 8 September 2026. Version~1, submitted on
23 September 2026, states that equivalence as its Theorem~1.4.
\end{remark}
\begin{proof}
\begin{enumerate}[label=\arabic*.]
\item Use the same oriented triangle, ordered integral basis
$(\widehat\gamma,\widehat u,\widehat w,\widehat\delta)$,
finite twists, and untranslated unit-square cusp fan as version~1.
Its chosen function $\beta_0$ and our $\beta^{\rm part}$ obey the
same affine transformation laws, and both $\beta_0+\tau$ and
$\beta^{\rm part}+\tau$ extend over the cusp. Their difference
descends to a holomorphic function on $\mathbb P^1$, hence is the
constant $\delta=\beta^{\rm part}-\beta_0$. The exact parameter
map and inverse are $a=c+\delta$ and $c=a-\delta$: substitution
gives $\Omega_{c+\delta}=\Pi_c$ in the stated ordered basis.
Writing $M_0=\sup(\Im\beta_0-6(\Im\mu)^2/\Im\tau)$,
our corresponding bound is $M_{\rm part}=M_0+\Im\delta$, so this
map also takes our admissible half-plane bijectively to version~1's.
The ordinary quotients, finite logarithmic maps, and cusp charts
then identify by their identical marked periods and specified
gluing data. In particular $a'-a=c'-c$. Its Theorem~1.4 therefore
gives the displayed exact-period equivalence and the full
automorphism group in this standard-data scope.
\item Theorem~6.12 of version~1 gives the tangent-cohomology tuple,
and Proposition~6.16 gives the smooth reduced Kuranishi germ and
the locally miniversal parameter family. These are the
four assertions compared with \cref{night:P5.3,night:P5.4,night:P5.5,night:P5.6}
under that same parameter and marking.
\item The second version retains the period construction and the
Oka conclusions. Its displayed section list and its mathematical
statements omit the first version's logarithmic-deformation and
classification sections and the four assertions just listed.
This comparison records the contents of the two specified versions;
it assigns no reason for the change.
\end{enumerate}
\end{proof}
